\documentclass[letterpaper,twocolumn,openany]{dndbook}

\usepackage[english]{babel}
\usepackage[utf8]{inputenc}

\title{The Master Herbalist's Compendium}
\author{Meadowleaf the Green}
\date{Compiled in the Year of the Wandering Wyrm}

\begin{document}

\frontmatter

\maketitle

\tableofcontents

\mainmatter

\chapter{Introduction to Herbalism}

\begin{DndReadAloud}
    Nature provides all the remedies we need, if only we know where to look and how to prepare them. This compendium represents a lifetime of gathering, experimentation, and occasionally, painful mistakes.

    Within these pages, you will find the secrets of poultice-making, antidote brewing, and the identification of both helpful and harmful plants.

    \hfill — Meadowleaf the Green, Druid of the Emerald Grove
\end{DndReadAloud}

\section{What is Herbalism?}

Herbalism is the art and science of using plants for medicinal, alchemical, and practical purposes. A skilled herbalist can:

\begin{itemize}
    \item Identify beneficial and poisonous plants
    \item Create healing poultices and salves
    \item Brew antidotes and antitoxins
    \item Craft herbal remedies for common ailments
    \item Prepare ingredients for alchemical formulas
\end{itemize}

\begin{DndComment}{Proficiency Required}
    To use an herbalism kit, you must be proficient with the tool. Herbalism kit proficiency can be gained through:
    \begin{itemize}
        \item Background (Hermit, Outlander)
        \item Class features (some Druids, Rangers)
        \item Downtime training (10 workweeks, 250 gp)
    \end{itemize}
\end{DndComment}

\chapter{The Herbalism Kit}

\section{Components}

A standard herbalism kit contains:

\begin{DndTable}[header=Kit Contents]{lX}
    \textbf{Item} & \textbf{Purpose} \\
    Pouches & Storing harvested herbs \\
    Mortar \& pestle & Grinding ingredients \\
    Glass vials & Storing liquids and extracts \\
    Linen strips & Creating poultices \\
    Knife & Harvesting and cutting \\
    Reference book & Plant identification \\
\end{DndTable}

\textbf{Cost:} 5 gp \\
\textbf{Weight:} 3 lbs.

\chapter{Common Herbs}

\section{Yarrow (Achillea millefolium)}

\begin{DndSidebar}{Yarrow Quick Reference}
    \textbf{Habitat:} Grasslands, meadows

    \textbf{Season:} Summer

    \textbf{Rarity:} Common

    \textbf{Harvesting DC:} 10

    \textbf{Primary Use:} Wound healing, fever reduction
\end{DndSidebar}

\subsection{Identification}

\begin{itemize}
    \item Feathery, fern-like leaves
    \item Flat-topped clusters of tiny white or pink flowers
    \item Strong, distinctive aroma
    \item Grows 1-3 feet tall
\end{itemize}

\subsection{Medicinal Properties}

Yarrow has been used for centuries to:

\begin{itemize}
    \item Stop bleeding (apply crushed leaves directly to wounds)
    \item Reduce fever (brew tea from dried flowers)
    \item Ease digestive discomfort (mild tea)
    \item Repel insects (hang dried bundles)
\end{itemize}

\subsection{Preparation Methods}

\textbf{Wound Poultice:}
\begin{enumerate}
    \item Crush fresh yarrow leaves in mortar
    \item Add small amount of water to create paste
    \item Apply directly to minor wounds
    \item Bind with clean linen
    \item Change daily until healed
\end{enumerate}

\textbf{Game Effect:} Stabilizes dying creature without expending hit dice (requires herbalism kit and 1 minute)

\section{Feverfew (Tanacetum parthenium)}

\begin{DndSidebar}{Feverfew Quick Reference}
    \textbf{Habitat:} Gardens, roadsides

    \textbf{Season:} Spring to early autumn

    \textbf{Rarity:} Common

    \textbf{Harvesting DC:} 12

    \textbf{Primary Use:} Headache relief, anti-inflammatory
\end{DndSidebar}

\subsection{Identification}

\begin{itemize}
    \item Daisy-like flowers with white petals and yellow centers
    \item Strongly scented, bitter-tasting leaves
    \item Grows in bushy clumps, 1-2 feet high
    \item Flowers from June to August
\end{itemize}

\subsection{Uses}

\begin{DndTable}[header=Feverfew Applications]{lX}
    \textbf{Ailment} & \textbf{Preparation} \\
    Headaches & Chew 2-3 fresh leaves daily \\
    Inflammation & Tea from dried flowers \\
    Fever & Cold infusion (overnight steep) \\
    Joint pain & Poultice of crushed leaves \\
\end{DndTable}

\textbf{Warning:} Some individuals develop mouth sores from chewing fresh leaves. Discontinue if irritation occurs.

\section{Comfrey (Symphytum officinale)}

\begin{DndComment}{The Boneset Herb}
    Comfrey is sometimes called "knitbone" for its remarkable ability to speed the healing of broken bones and torn tissue. However, it should NEVER be taken internally in large quantities.
\end{DndComment}

\subsection{Identification}

\begin{itemize}
    \item Large, hairy leaves
    \item Purple, pink, or white bell-shaped flowers
    \item Deep taproot (important for harvesting)
    \item Grows near water sources
    \item 2-4 feet tall when flowering
\end{itemize}

\subsection{Medicinal Uses}

\textbf{External use only:}

\begin{itemize}
    \item Broken bones (poultice on external area)
    \item Sprains and strains
    \item Bruises and contusions
    \item Minor cuts (speeds healing)
\end{itemize}

\textbf{Preparation:}

\begin{enumerate}
    \item Harvest fresh root (or use dried)
    \item Clean thoroughly and chop finely
    \item Boil in water for 10 minutes to create decoction
    \item Soak linen in warm decoction
    \item Apply to affected area
    \item Wrap and leave for several hours
\end{enumerate}

\textbf{Game Effect:} During a short rest, allows an additional hit die to be spent for healing (requires herbalism kit and comfrey)

\chapter{Rare and Magical Herbs}

\section{Silverleaf (Rare)}

\begin{DndSidebar}{Silverleaf}
    \textbf{Habitat:} Moonlit glades, elven territories

    \textbf{Rarity:} Rare

    \textbf{Harvesting DC:} 18

    \textbf{Value:} 50 gp per leaf

    \textbf{Special:} Only visible under moonlight
\end{DndSidebar}

\subsection{Description}

Silverleaf is a magical plant that appears as an ordinary weed by day but reveals shimmering silver leaves under moonlight.

\subsection{Properties}

\begin{itemize}
    \item Antitoxin ingredient (3 leaves required)
    \item Delays poison effects by 1 hour when chewed
    \item Essential for crafting potions of resistance (poison)
    \item Highly valued by alchemists
\end{itemize}

\textbf{Harvesting:} Must be gathered during a full moon. The plant withers by dawn if not properly preserved in silver-lined containers.

\section{Dragon's Breath Pepper (Very Rare)}

\begin{DndReadAloud}
    I once saw a halfling eat an entire Dragon's Breath Pepper on a dare. He breathed actual fire for the next minute—small flames, mind you, but real fire nonetheless. He won the dare but lost his sense of taste for a month.
\end{DndReadAloud}

\subsection{Properties}

\begin{itemize}
    \item Extreme heat (literally magical)
    \item Component for fire resistance potions
    \item Can be weaponized (thrown, causes 1d4 fire damage and blindness)
    \item Culinary use (extreme spice, 1/100th of a pepper flavors an entire pot)
\end{itemize}

\textbf{Warning:} Handle with leather gloves. Wash hands thoroughly after contact. Keep away from eyes, nose, and sensitive areas.

\chapter{Crafting Recipes}

\section{Basic Healing Poultice}

\textbf{Time:} 10 minutes \\
\textbf{DC:} 15 (Wisdom check using herbalism kit) \\
\textbf{Cost:} 5 gp in materials

\textbf{Ingredients:}
\begin{itemize}
    \item 3 sprigs of fresh yarrow
    \item 2 comfrey leaves
    \item 1 dose of honey (preservative and antibacterial)
    \item Clean linen bandage
\end{itemize}

\textbf{Method:}
\begin{enumerate}
    \item Crush yarrow and comfrey together
    \item Add honey to create thick paste
    \item Spread on linen bandage
    \item Apply to wound
\end{enumerate}

\textbf{Effect:} When applied during a short rest, grants advantage on next hit die rolled for healing.

\section{Antitoxin}

\textbf{Time:} 1 hour \\
\textbf{DC:} 20 (Wisdom check using herbalism kit) \\
\textbf{Cost:} 50 gp in materials

\textbf{Ingredients:}
\begin{itemize}
    \item 3 silverleaf leaves
    \item 1 crushed feverfew flower
    \item Purified water
    \item Glass vial
\end{itemize}

\textbf{Method:}
\begin{enumerate}
    \item Boil purified water under moonlight
    \item Add silverleaf (water turns pale blue)
    \item Steep for 30 minutes
    \item Add crushed feverfew
    \item Strain into vial
    \item Seal with wax
\end{enumerate}

\textbf{Effect:} Functions as standard antitoxin (grants advantage on saving throws vs. poison for 1 hour)

\chapter{Foraging and Harvesting}

\section{Finding Herbs}

To forage for specific herbs:

\begin{enumerate}
    \item Spend 1 hour searching appropriate terrain
    \item Make Wisdom (Survival) or Intelligence (Nature) check
    \item Consult the Foraging Table below
\end{enumerate}

\begin{DndTable}[header=Foraging Success]{cX}
    \textbf{Check Result} & \textbf{Herbs Found} \\
    10-14 & 1d4 common herbs \\
    15-19 & 1d6 common herbs \\
    20-24 & 1d6 common + 1d4 uncommon herbs \\
    25+ & 1d6 common, 1d4 uncommon, 1 rare herb \\
\end{DndTable}

\section{Preservation}

Proper storage is essential:

\begin{DndTable}[header=Herb Storage]{lX}
    \textbf{Type} & \textbf{Method} \\
    Fresh herbs & Use within 3 days, or dry \\
    Dried herbs & Airtight container, dark location \\
    Roots & Bury in sand or dry and store \\
    Flowers & Hang upside-down to dry \\
    Seeds & Store in paper envelopes \\
\end{DndTable}

\textbf{Shelf Life:}
\begin{itemize}
    \item Dried herbs: 1 year
    \item Dried roots: 2 years
    \item Seeds: 3 years
    \item Fresh: 3 days
\end{itemize}

\chapter{Safety and Ethics}

\section{Ethical Harvesting}

\begin{DndComment}{The Herbalist's Oath}
    "Take only what you need. Leave enough for regrowth. Thank the plant for its sacrifice. Restore the earth when you depart."
\end{DndComment}

\textbf{Sustainable practices:}

\begin{enumerate}
    \item Never take more than 1/3 of available plants
    \item Harvest from multiple locations
    \item Avoid rare species unless truly necessary
    \item Replant or scatter seeds when possible
    \item Fill in holes after root harvesting
    \item Leave offerings to nature spirits (optional but respectful)
\end{enumerate}

\section{Poisonous Plants to Avoid}

\begin{DndTable}[header=Dangerous Plants]{lX}
    \textbf{Plant} & \textbf{Danger} \\
    Wolfsbane & Extremely toxic, causes paralysis \\
    Nightshade & Deadly poison, resembles edible berries \\
    Hemlock & Fatal if ingested, smells like parsnips \\
    Deathcap Mushroom & No antidote, causes liver failure \\
\end{DndTable}

\textbf{Rule:} If you're not 100\% certain of identification, do not harvest or consume.

\chapter{Conclusion}

\begin{DndReadAloud}
    You now hold the knowledge to become a skilled herbalist. But remember: knowledge must be tempered with wisdom and respect for nature.

    The plants give freely of their healing power. In return, we must be their stewards and protectors.

    Go forth and heal. May your hands always find the right herb at the right moment.

    \hfill — Meadowleaf the Green
\end{DndReadAloud}

\section{Herbalist's Checklist}

\begin{itemize}
    \item[$\square$] Can identify 10+ common medicinal herbs
    \item[$\square$] Know proper harvesting techniques
    \item[$\square$] Can create basic healing poultice
    \item[$\square$] Understand ethical foraging practices
    \item[$\square$] Recognize dangerous/poisonous plants
    \item[$\square$] Know proper storage methods
    \item[$\square$] Can brew antitoxin
    \item[$\square$] Respect nature's balance
\end{itemize}

\vfill

\begin{center}
    \textit{In herbs, we find nature's pharmacy.}
\end{center}

\end{document}
